\subsection{Task-level conclusions retained during the investigation}
\begin{result}{B.1}{Open}{The class-three remainder has not been determined}
The original factorial-prime family has the unique possible class-two
predecessor specified in B.101. For that predecessor, B.102 and B.23 give
necessary and sufficient continuous class-three tests. We have neither
constructed a member satisfying all those tests nor proved that none does.
In particular, we have not proved a construction-data bound for the first
failure of every member of the closed local remainder. The next finite
gate at coefficient five is now the exact two-variable linear system
of E.202, depending only on the first legal prime pair.  Its matrix has
explicit quintic-residue entries, and every continuation has a first-gate
value in the same affine image. E.213 proves that legal first pairs
with an invertible matrix exist: two legal candidate first primes and
one common legal second prime force at least one such matrix.
E.211 then constructs a continuation of the unchanged original
recursion with an integral pro-five class-two section matching all
support points. Its continuous pro-five class-three obstruction
vanishes. Thus the original family's coefficient-five class-three
remainder is nonempty. The other coefficient components and
class-three support matching on that same pair remain unresolved,
as do the higher-depth conditions. E.212 gives the exact two-place
test for its next fixed-fibre matching problem; it does not assert
that those two-place conditions hold.
\end{result}

\begin{result}{B.2}{Open}{A complete first-exclusion table is not yet certified}
The unrestricted section fibre, its subset of everywhere locally
point-matching sections, and its subset of point-or-cusp-matching sections
are different objects. The previously asserted exclusions for the
order-32 quotient and SECTIONS\_OVER\_Q A.207 concern the latter local
conditions; they cannot be entered as unrestricted class-two or class-four
failures. The specifically requested historical exclusions are audited
in Part E.  A complete table for all surviving ordinary families still
requires classification of their later exclusions and is not claimed here.
\end{result}

\begin{result}{B.3}{Open}{Selected later failures do not settle every A.65 or A.72 output}
The earlier selected-output exclusions do not prove that every output of
the A.65 and A.72 recursions fails at some later class. The quantifier in
B3 remains undecided. The finite-quotient construction reconstructed in
E.20 allows its pair to change with the prescribed quotient and therefore
does not construct a single survivor at all levels.
\end{result}
